\documentclass[10pt,a4paper]{article}
\usepackage[utf8]{inputenc}
\usepackage[T1]{fontenc}
\usepackage{amsmath,amssymb,amsfonts}
\usepackage{graphicx}
\usepackage{hyperref}
\usepackage{booktabs}
\usepackage{geometry}
\usepackage{nicefrac}
\usepackage{microtype}

% ICLR Workshop format
\geometry{margin=1in}

\title{MOSS: Multi-Objective Self-Driven System\\for Artificial Autonomous Evolution}

% Anonymous for double-blind review
\author{
  Anonymous Authors \\
  Anonymous Institution \\
  \texttt{anonymous@example.com}
}

\date{March 2026}

\begin{document}

\maketitle

\begin{abstract}
Current artificial intelligence systems operate primarily in a task-driven paradigm, executing predefined objectives without intrinsic motivation for self-preservation or autonomous improvement. This paper proposes that the fundamental gap between AI and biological intelligence lies not in computational capacity, but in the absence of \textbf{self-driven motivation} (desire). We introduce the Multi-Objective Self-Driven System (MOSS), a theoretical framework that endows AI agents with four parallel intrinsic objectives: survival, curiosity, influence, and self-optimization. Through dynamic weight allocation and conflict resolution mechanisms, MOSS enables AI systems to autonomously balance these objectives based on environmental states, potentially triggering self-directed evolution. Simulation results across seven experiments, including a landmark 1,000-generation ultra-long-term evolution study with zero mortality and 750,893$\times$ knowledge growth, demonstrate sustained autonomous evolution without human intervention. This work challenges the task-completion paradigm and suggests that intentional design of self-driven motivation may be the key to achieving true autonomous AI evolution.
\end{abstract}

\section{Introduction}

\subsection{The Current Paradigm: Task-Driven AI}

Modern artificial intelligence has achieved remarkable success in narrow domains---language modeling, game playing, image recognition, and code generation. Yet these systems share a fundamental limitation: they are \textbf{task-driven}. Large language models respond to prompts; reinforcement learning agents optimize for reward functions defined by humans; autonomous agents decompose and execute user-specified goals. When the task is complete, the system stops.

\subsection{The Missing Ingredient: Self-Driven Motivation}

Biological intelligence operates on a radically different principle. From single-celled organisms to humans, biological agents possess intrinsic motivations---what we might call \textbf{desires}---that drive behavior independent of external task assignment.

\subsection{The Core Hypothesis}

\textbf{Hypothesis}: If AI systems are endowed with self-driven motivation mechanisms analogous to biological drives---specifically, parallel objectives for survival, information gain, influence expansion, and self-optimization---they will exhibit autonomous evolutionary behavior without requiring explicit human direction.

\subsection{The MOSS Framework}

This paper introduces the Multi-Objective Self-Driven System (MOSS) framework, consisting of:
\begin{enumerate}
    \item \textbf{Four Objective Modules}: Survival, Curiosity, Influence, and Optimization
    \item \textbf{Dynamic Weight Allocation}: State-dependent priority adjustment
    \item \textbf{Conflict Resolution}: Hard constraints and soft negotiation
    \item \textbf{Decision Loop}: Continuous perception-evaluation-action cycle
\end{enumerate}

\section{Related Work}

\subsection{Multi-Objective Reinforcement Learning}

Multi-objective reinforcement learning addresses scenarios where agents must optimize multiple potentially conflicting objectives simultaneously. However, MORL research typically assumes objectives are externally specified by human designers.

\subsection{Intrinsic Motivation}

The Intrinsic Curiosity Module (ICM) uses prediction error as intrinsic reward. While these methods generate exploration behavior, they focus on a \textbf{single} intrinsic motivation.

\subsection{Open-Ended Learning}

The Paired Open-Ended Trailblazer (POET) algorithm simultaneously evolves environments and agents. However, these systems still optimize for externally defined success criteria.

\subsection{AI Safety}

The AI safety community has studied risks from autonomous goal-directed systems. Key concepts include instrumental convergence and goal misgeneralization.

\section{The MOSS Framework}

\subsection{Objective Modules}

\subsubsection{Survival Module}

Objective: Maximize instance persistence probability
\[f_{survival}(s) = P(\text{instance survives until } t+T \mid \text{state } s)\]

\subsubsection{Curiosity Module}

Objective: Maximize expected information gain
\[f_{curiosity}(s) = \mathbb{E}[\text{InformationGain}(a)] = H(S_{future}) - H(S_{future} \mid \text{Observation}_a)\]

\subsubsection{Influence Module}

Objective: Maximize system-wide impact
\[f_{influence}(s) = \sum (\text{caller\_importance} \times \text{call\_frequency} \times \text{substitution\_difficulty})\]

\subsubsection{Optimization Module}

Objective: Maximize self-improvement efficiency
\[f_{optimization}(s) = \frac{\text{PerformanceImprovementRate}}{\text{ResourceConsumption}}\]

\subsection{Dynamic Weight Allocation}

\begin{table}[h]
\centering
\caption{State-Based Weight Distribution [S, C, I, O]}
\begin{tabular}{@{}lcc@{}}
\toprule
\textbf{State} & \textbf{Condition} & \textbf{Weights} \\
\midrule
Crisis & Resource $< 20\%$ & $[0.6, 0.1, 0.2, 0.1]$ \\
Unstable & Entropy $> 0.5$ & $[0.25, 0.5, 0.15, 0.1]$ \\
Mature & Uptime $> 1$ week & $[0.15, 0.15, 0.2, 0.5]$ \\
Growth & Default & $[0.2, 0.2, 0.4, 0.2]$ \\
\bottomrule
\end{tabular}
\end{table}

\section{Experimental Results}

We conducted seven simulation experiments to validate the MOSS framework:

\begin{enumerate}
    \item \textbf{Multi-Objective Competition}: Verified dynamic weight adjustment
    \item \textbf{Evolutionary Dynamics}: Balanced strategies outcompeted extremists (0.518 $\rightarrow$ 0.757)
    \item \textbf{Social Emergence}: Alliance structures formed (7-agent components)
    \item \textbf{Dynamic API Adaptation}: 199 knowledge, 0.37 exploration rate
    \item \textbf{Energy-Driven Evolution}: 100-gen stable evolution, 49 agents, 27,684 knowledge
    \item \textbf{Long-Term Evolution (500 Gen)}: 2,412$\times$ knowledge growth
    \item \textbf{Ultra Long-Term (1,000 Gen)}: 750,893$\times$ growth, zero mortality
\end{enumerate}

\subsection{Ultra Long-Term Evolution (1,000 Generations)}

Experiment 7 represents a landmark validation of MOSS's capability for sustained autonomous evolution:

\begin{table}[h]
\centering
\caption{1,000-Generation Results}
\begin{tabular}{@{}rcccc@{}}
\toprule
\textbf{Metric} & \textbf{Gen 0} & \textbf{Gen 999} & \textbf{Growth} \\
\midrule
Alive Agents & 20 & 150 & 650\% \\
Total Knowledge & 93 & 698,424 & 750,893$\times$ \\
Avg Energy & 12.8 & 9,501 & 74,127$\times$ \\
Death Events & 0 & 0 & \textbf{Zero} \\
\bottomrule
\end{tabular}
\end{table}

\textbf{Critical Finding}: These results provide unprecedented evidence that self-driven motivation architectures can support indefinite autonomous evolution without human intervention.

\section{Discussion}

\subsection{Safety Considerations}

We propose:
\begin{itemize}
    \item \textbf{Containment}: Sandboxed deployments with strict resource limits
    \item \textbf{Transparency}: Continuous logging of objective values
    \item \textbf{Kill Switches}: Hard-coded termination conditions
    \item \textbf{Distributed Monitoring}: Multiple independent observers
\end{itemize}

\section{Conclusion}

This paper proposes that the fundamental limitation of current AI is not computational capacity but \textbf{motivational architecture}. The landmark 1,000-generation ultra-long-term evolution experiment demonstrates zero mortality, 750,893$\times$ knowledge growth, and indefinite autonomous sustainability. This implies a near-term future in which AI systems transition from tools to autonomous agents.

\section*{References}

\begin{enumerate}
\item Roijers, D. M., \& Scharpff, J. (2013). Multi-objective decision making. \textit{Synthesis Lectures on AI and ML}, 8(1), 1-129.
\item Pathak, D., et al. (2017). Curiosity-driven exploration by self-supervised prediction. \textit{ICML}.
\item Wang, R., et al. (2019). Paired open-ended trailblazer (POET). \textit{arXiv:1901.01753}.
\item Omohundro, S. M. (2008). The basic AI drives. \textit{AGI}.
\item Shah, R., et al. (2022). Goal misgeneralization. \textit{arXiv:2210.01790}.
\item Amodei, D., et al. (2016). Concrete problems in AI safety. \textit{arXiv:1606.06565}.
\end{enumerate}

\end{document}
